\section{Background and Related Work}

Ransomware detection remains a critical area of research, driven by the continuous evolution of extortion tactics. The security community has explored a spectrum of methodologies, ranging from signature-based matching to advanced machine learning models. This section reviews principal approaches and situates ADRDS within the existing body of work.

\subsection{Evolution of Detection Techniques}

\textbf{Signature-Based Methods.}
Traditional antivirus solutions rely on maintaining databases of known malware signatures. While effective against established threats, this paradigm is insufficient against polymorphic ransomware that dynamically alters its code, or zero-day variants lacking prior cataloging. Consequently, signature-based reliance is increasingly untenable.

\textbf{Behavior-Based Detection.}
Behavioral analysis monitors execution patterns rather than static code attributes. By tracking file system activity, registry modifications, process behaviors, and system call sequences, detection systems can identify emerging attacks. Pioneering work by Continella et al. \cite{continella2016shieldfs} and Scaife et al. \cite{scaife2016cryptolock} demonstrated that monitoring abnormal file access and encryption rates facilitates early detection. However, distinguishing malicious activity from legitimate operations performed by backup utilities or compression tools remains a challenge, often leading to false positives.

\textbf{Machine Learning Approaches.}
To enhance detection precision, recent research has leveraged machine learning. Supervised classifiers—such as Random Forests, Support Vector Machines, and gradient-boosted trees—trained on labeled datasets can identify subtle patterns compliant with known attack vectors. Deep learning models, including recurrent neural networks and autoencoders, have also shown efficacy in capturing temporal behaviors and anomalies. Studies by Sgandurra et al. \cite{sgandurra2016automated} and Vinayakumar et al. \cite{vinayakumar2019deep} report promising results; however, many existing systems operate offline, exhibit high latency, or function as opaque models lacking interpretability.

\textbf{Hybrid Systems.}
Hybrid detection systems integrate multiple methodologies—static and dynamic analysis, alongside supervised and unsupervised learning—to improve resilience. Research by Kolosnjaji et al. \cite{kolosnjaji2018deep} and Pascanu et al. \cite{pascanu2013malware} indicates that hybrid models offer superior generalization and noise tolerance. Nevertheless, these systems often incur significant computational overhead and typically lack interactive capabilities suitable for educational or exploratory use.

\subsection{Interactive and Demonstration Platforms}

Despite advances in detection accuracy, limited tools prioritize hands-on exploration. Most platforms operate as automated backend services with minimal visibility into decision-making processes. existing educational sandboxes, such as CyberSecLab \cite{lee2022cyberseclab}, enable experimentation with general security concepts but often lack specialized ransomware behavioral taxonomies, real-time machine learning integration, and transparent logic visualization.

\subsection{System Positioning}

ADRDS addresses this gap by providing an interactive environment. Unlike "black-box" systems, ADRDS enables users to trigger simulated ransomware behaviors, adjust parameters, and observe the detection engine's response in real time. The system combines safe behavioral simulation with a hybrid ML detector and an interactive dashboard. This approach provides both rigorous detection capabilities and the necessary transparency to understand alert triggers, making ADRDS valuable for advanced security research and instructional applications.

\subsection{Foundational Research}

The architecture and evaluation of ADRDS are informed by several key studies. Kharraz et al. \cite{kharraz2015cutting} provide analysis of the fundamental mechanisms of ransomware attacks. Continella et al. \cite{continella2016shieldfs} propose ShieldFS, a self-healing filesystem that influenced the behavioral monitoring approach used here. Scaife et al. \cite{scaife2016cryptolock} introduce CryptoLock, focusing on preventing data loss. Sgandurra et al. \cite{sgandurra2016automated} provide a framework for automated dynamic analysis. Vinayakumar et al. \cite{vinayakumar2019deep} survey deep learning approaches, while Kolosnjaji et al. \cite{kolosnjaji2018deep} explore hybrid analysis using system call sequences. Pascanu et al. \cite{pascanu2013malware} demonstrate the utility of Recurrent Neural Networks (RNNs) in malware classification, and Lee et al. \cite{lee2022cyberseclab} present CyberSecLab, which inspired the educational and interactive focus of this work.
